
\section{Metodologia}

\begin{frame}
\frametitle{Algoritmo de SVD}

\begin{itemize}
    \item Para una matrix $X \in \mathbb{R^{T\times D}}$, con $T$ filas y $D$ columnas, se busca descomponer $X$ de la siguiente forma:
    \begin{equation}
        X = U\Sigma V^T \Leftrightarrow XV = U\Sigma
    \end{equation}
    \item Donde puede identificarse una nueva matriz $W = XV$. El objetivo del algoritmo es generar $W$ y $V$ a partir de $X$.
    \item Para ello, se pueden rotar iterativamente todos los posibles pares de columnas en $X$ hasta que los valores no diagonales de la matriz $X^TX$ converjan
    a un cierto valor de tolerancia $\epsilon$ ($\mathcal{O}(D^2)$).
\end{itemize}

\end{frame}

\begin{frame}
\frametitle{Algoritmo de SVD}
\begin{itemize}
    \item El proceso de ortogonalización de las columnas se realiza mediante \textbf{rotaciones de Givens}, buscando una matriz de rotación $G$, tal que el producto de las columnas con $G$ sean ortogonales.
    \begin{equation}
        G = \left[ 
            \begin{array}{cc}
                \cos \theta & -\sin \theta \\
                \sin \theta & \cos \theta.
            \end{array}
         \right]
    \end{equation}
    \item Bajo ciertas estrategias matemáticas, estos valores pueden obtenerse definiendo un parámetro especial $\tau$, definido como:
    \begin{equation}
    \tau = \frac{x_j^T\cdot x_j - x_i^T \cdot x_j}{2 (x_i^T \cdot x_j)}
    \end{equation}
    \item Definido por el productos punto entre las columnas a rotar ($x_i, x_j$) y los productos punto consigo mismas. ($\mathcal{O}(T)$)
\end{itemize}
\end{frame}

\begin{frame}
\frametitle{Algoritmo de SVD}
\begin{itemize}
    \item Previo a este proceso iterativo se inicializa una matriz identidad $V$ de tamaño $D \times D$.
    \item Se aplica la rotación obtenida a las columnas de $X$ usando la matriz $G$, y adicionalmente se ''vuelcan'' acumulativamente estas rotaciones en la columnas correspondientes de $V$. ($\mathcal{O} (D)$).
    \item Una vez que se alcanza la convergencia al valor $\epsilon$ en todas las columnas, la matriz $X$ pasa a ser $W$, tal que $W = XV$. En base a $W$ pueden obtenerse tanto $\Sigma$ como $U$.
\end{itemize}
\end{frame}

\begin{frame}
\frametitle{Complejidad}
\begin{itemize}
    \item Considerando las complejidades anteriores, asumiendo un proceso secuencial con iteraciones anidadas, la complejidad de este algoritmo sería de
    $\mathcal{O}(n(D^2(T+D)))$, con $n$ el número de iteraciones necesarias para alcanzar la convergencia.
\end{itemize}
\end{frame}

\begin{frame}
\frametitle{Paralelización}
\begin{itemize}
    \item Dentro del algoritmo, el proceso de ortogonalización iterativo por cada par de columnas puede paralelizarse, definiendo un algoritmo de ''barrido''
    por todos los pares de columnas y definiéndolo como un kernel para una implementación en CUDA.
\end{itemize}
\end{frame}

\begin{frame}
\frametitle{Pseudocódigo}
\begin{algorithm}[H]
    \tiny
    \caption{CUDA based One-Sided Jacobi implementation}
    \begin{algorithmic}[1]
        \State \textbf{Inputs:}\par $\mathbf{X} \gets \text{data matrix}\, (T \times D)$; \par
        $\epsilon \gets \text{tolerance}$; \par
        $\texttt{max\_sweeps} \gets \text{maximum amount of sweeps}$;
        \State \textbf{Initializations:}\par
        $I \gets \{0,\dots,D-1\},\, \text{indexes of}\, \mathbf{X}$;\par
        $\mathbf{V} \gets \text{identity matrix of size }D\times D$;\par
        $\texttt{any\_rots} \gets \text{loop condition}$;\par
        \State \text{Initialize }$\mathbf{X}, \mathbf{V}$\text{ to }\texttt{Device};
        \State \texttt{any\_rots\_host }$\gets$\texttt{ True};
        \State \texttt{sweeps }$\gets 0,$\text{ sweep counter};
        \While {(\texttt{any\_rots\_host} \textbf{and} \texttt{sweeps }$<$\texttt{ max\_sweeps})}
            \State \texttt{any\_rots\_host }$\gets$\texttt{ False};
            \State \texttt{any\_rots\_device }$\gets$\texttt{ any\_rots\_host},\text{ copy host value to device};
            \For {($r$\textbf{ from }$0$\textbf{ to }$D-1$)}
                \State \texttt{Sweep }$\lll D/2\text{ blocks, }1\text{ thread}\ggg(\mathbf{X},\mathbf{V},I_D,\texttt{any\_rots\_device})$\text{; execute device kernel}
                \State \textbf{synchronize threads};
                \State \texttt{circular\_shift}($I$),\, \text{shuffle indexes using }\textit{circular shift};
            \EndFor
            \State \texttt{any\_rots\_host }$\gets$\texttt{ any\_rots\_device};
            \State \texttt{sweeps++};
        \EndWhile
    \end{algorithmic}
\end{algorithm}
\end{frame}

\begin{frame}
\frametitle{Kernel propuesto}
\begin{algorithm}[H]
    \scriptsize
    \caption{\texttt{Sweep} kernel}
    \begin{algorithmic}[1]
        \State \textbf{Sweep}($\mathbf{X},\mathbf{V},I_D,\texttt{any\_rots\_device}$)
        \For{\textbf{each }\text{pair(i,j) }\textbf{in }$I_D$}
            \State $a \gets \mathbf{X}[\text{i}] \cdot \mathbf{X}[\text{i}]$;
            \State $b \gets \mathbf{X}[\text{j}] \cdot \mathbf{X}[\text{j}]$;
            \State $c \gets \mathbf{X}[\text{i}] \cdot \mathbf{X}[\text{j}]$;
            \State \texttt{needs\_rot }$\gets(|c|/\sqrt{a\cdot b}) > \epsilon$\text{, rotation condition};
            \If{\texttt{needs\_rot}}
                \State $\tau, t,$\texttt{cos\_theta},\texttt{ sin\_theta }$\gets $\texttt{ given\_rot}($\mathbf{X}[\text{i}], \mathbf{X}[\text{j}]$);
                \State \texttt{update}($\mathbf{X}[\text{i}], \mathbf{X}[\text{j}]$);
                \State \texttt{update}($\mathbf{V}[\text{i}], \mathbf{V}[\text{j}]$);
                \State \texttt{any\_rots\_device} = \texttt{any\_rots\_device} \textbf{or} 1;
            \EndIf
        \EndFor
    \end{algorithmic}
\end{algorithm}
\end{frame}

\begin{frame}
\frametitle{Algoritmo paralelo}
\begin{itemize}
    \item El algoritmo utiliza una serie de ''barridos'' utilizando un arreglo de índices, realizando un proceso de \textit{circular shift}, asegurando que la repetición de pares sea mínima.
    \item Cada hilo realiza el proceso de rotación de manera independiente, comprobando si el par a rotar ya es ortogonal en sí mismo.
    \item Esta componente de paralelismo reduciría la complejidad del barrido de $\mathcal{O} (D^2)$ a $\mathcal{O}(D)$, dado que solo tienen que realizarse $D-1$ barridos como máximo.
\end{itemize}
\end{frame}


